\section{Known PPM Sequence Delay MLE}
\label{sec:Known-PPM-sequence-Delay-Estimation}

In this section, we finally look at the problem when the transmitted signal \(t\mapsto s(t)\) is a known PPM signal with inter-symbol guard times (ISGTs). We assume that every symbol is composed by \(n\) slots, where messages of length \(\log_2n\) are encoded, followed by \(P\) slots not containing any signal pulses. The length of a slot is then \(T_s=T_{\text{symbol}}/(n+P)\). As in~\Cref{sec:Repeated-Pulse-Delay-Estimation} we assume to be looking at a total of \(M\) symbols. As in~\eqref{eq:aparently} we have:
\begin{equation}
y_r[k]=\sum_{m=0}^{M-1}x_r[k+mQ] \quad k\in\{0,\ldots,Q-1\}.
\label{eq:aparently2}
\end{equation}
where \(Q=n+P\).

\noindent
If we assume the signal slots to be deterministically uniform (i.e. the same number of pulses appear in each signal slot over the duration of the \(M\) symbols), then there are \(M/n\) aggregated pulses in \(y_r[k]\) for \(k\in\{0,\ldots,n-1\}\).

\noindent
With this set up the problem is almost identical to~\eqref{eq:position} with only a difference in the constant multiplying the periodic signal:
\begin{equation}
y_r[k]\sim
\operatorname{Poi}
\left(\xi_{k,r}(\delta_r)\right),
\end{equation}
with
\begin{equation}
\begin{aligned}
\xi_{k,r}(\delta_r)
=  M\lambda_{b,r} T_s+\frac{M}{n}\lambda_{s,r} \int_{kT_s}^{(k+1)T_s}  \operatorname{rect}_{(nT_s)} \left((t-\delta_r) \,\text{mod}\, QT_s\right) \di t.
\end{aligned}
\end{equation}

\subsection{ML Estimation PPM Sequence}

\noindent
The result is almost exactly the same as in~\Cref{sub:MLestimation2} with a slight difference with the constants. When we then take the derivative with the log-likelihood over the interval \(\delta \in(kT_s,(k+1)T_s)\) we now get:
\begin{equation}
\begin{aligned}
\frac{\partial}{\partial \delta} \ell(\delta ; \mathbf{x})
&= \frac{\partial}{\partial \delta}\sum_{k=0}^{Q-1} y_k \log \left( \xi_k(\delta)\right)
=  \sum_{k=0}^{Q-1}y_k\frac{ \frac{\partial}{\partial \delta}\xi_k(\delta)}{\xi_k(\delta)}\\
&= y_{k+n} \frac{(M/n)\lambda_s}{\left((M/n)\lambda_s \delta-k (M/n)\lambda_s T_s\right)+M\lambda_b T_s}\\
&-y_k \frac{(M/n)\lambda_s}{\left(-(M/n)\lambda_s \delta+(k+1) (M/n)\lambda_s T_s\right)+M\lambda_b T_s}\\
&= \frac{y_{k+n}}{\left( \delta-k  T_s\right)+\frac{n\lambda_b}{\lambda_s} T_s}- \frac{y_k}{\left(- \delta+(k+1)  T_s\right)+\frac{n\lambda_b}{\lambda_s} T_s}.
\end{aligned}
\end{equation}

\noindent
From this we know that:
\begin{equation}
\begin{aligned}
\delta_{\mathrm{ML}, i}
&=T_s\left[\frac{y_{i+n}\left(i+1 +\frac{n\lambda_b}{\lambda_s} \right)+y_i\left(i -\frac{n\lambda_b}{\lambda_s} \right)}{y_i+y_{i+n}}\right],
\end{aligned} 
\end{equation}
and the final estimate is then found as the maximum between the interval estimates:
\begin{equation}
\delta_{\mathrm{ML}}=\underset{i}{\argmax} \; \ell\left(\delta_{\mathrm{ML}, i} ; \mathbf{y}\right) \quad i\in\{0, \ldots, Q-1\}.
\end{equation}
The Cramér-Rao Bound, the bias and the variance of the are the same as the ones found in~\Cref{sec:Single-Pulse-Delay-Estimation} and~\Cref{sec:Repeated-Pulse-Delay-Estimation} only with \(y_r[k]\) as defined in~\eqref{eq:aparently2} instead of \(x_r[k]\) and replacing \(\lambda_b\) with \(n\lambda_b\).